\chapter*{Chapter 3}
\markboth{Chapter 3}{3 Implementation - Phase 1: Sprints 1-3}
\addcontentsline{toc}{chapter}{3 Implementation - Phase 1: Sprints 1-3}

\textsc{\textbf{\Huge Implementation - Phase 1}} \\ \\

\stepcounter{chapter}

% ============================================================
\section*{Introduction}
\addcontentsline{toc}{section}{Introduction}

This chapter presents the implementation work completed during
the first phase of the project, covering Sprints 1 through 3.
This phase established the fundamental infrastructure that
the entire platform depends upon. Unlike the analysis and
architecture work of Sprint 0, these sprints focused on
delivering functional features that users can interact with
directly. The three main areas of focus were the user
authentication system, the product catalog and search
functionality, and the automated data scraping pipeline.
Each sprint built upon the previous one, creating a solid
foundation for all subsequent development work. The decisions
made during implementation prioritized code quality,
maintainability, and scalability to ensure the platform could
grow beyond its initial scope.

% ============================================================
\section*{Sprint 1: Authentication and User Management}
\addcontentsline{toc}{section}{Sprint 1: Authentication and User Management}

\subsection*{Sprint Overview}

Sprint 1 concentrated on building a robust authentication system
capable of handling multiple user types with different access
levels. The primary objective was to create a secure login and
registration process that could distinguish between five distinct
user roles: visitors, clients, merchants, suppliers, and
administrators. This sprint was critical because all other
features depend on correctly identifying and authorizing users.
Without a reliable authentication system, there would be no way
to personalize experiences, save preferences, or protect sensitive
functionality from unauthorized access.

The platform needed to support visitors who browse without
creating an account, clients who need personalized features like
favorites and price alerts, merchants who manage their product
offerings, suppliers who operate the data scraping tools, and
administrators who oversee the entire system. Each role required
access to different parts of the application, so the permission
system needed to be granular and reliable.

A two-step registration process was implemented to enhance
security and ensure data quality. Users first create their
account with basic information, then verify their phone number
through a one-time code sent by SMS. This verification step
serves multiple purposes: it confirms the phone number is valid,
enables future price alert notifications, and adds a layer of
protection against fake accounts.

The technology stack for this sprint centered on Laravel Fortify
for authentication logic, Laravel Sanctum for API token generation,
and Twilio for SMS delivery of verification codes. This
combination provided a mature, well-tested foundation that
integrated smoothly with the existing Laravel framework.

\subsection*{Implementation Details}

The user model was extended with a role field that determines
what features each user can access. Rather than creating separate
tables for each user type, a single users table holds all
accounts with a role column that can be set to client, merchant,
employee, supplier, or admin. This approach simplifies the
database structure and makes it easy to add new roles in the
future without modifying the schema.

\begin{table}[H]
\begin{center}
\begin{tabular}{|m{3cm}|m{6cm}|m{4cm}|}
\hline
\textbf{Model} & \textbf{Key Attributes} & \textbf{Purpose} \\ \hline
User & name, prename, email, phone, role & Core authentication entity \\ \hline
PhoneOtp & phone, code, expires\_at & Temporary verification codes \\ \hline
Client & user\_id, phone & Extended profile for clients \\ \hline
Merchant & user\_id, company\_name, is\_verified & Extended profile for merchants \\ \hline
Employee & user\_id, position & Platform operators \\ \hline
Supplier & user\_id, api\_key, merchant\_website\_id & Data provider access \\ \hline
\end{tabular}
\caption{Authentication Models}
\label{tab:auth_models}
\end{center}
\end{table}

The authentication flow follows a clear sequence. When a user
submits their registration details, the system generates a
verification code and stores it in the PhoneOtp table with an
expiration time. The code is sent via Twilio to the phone number
provided. The user then enters this code to complete registration.
Only after successful verification does the account become active.

During login, the system validates credentials and generates a
JWT token using Laravel Sanctum. This token must be included in
all subsequent API requests to authenticate the user. The token
expires after a configurable time period, requiring users to
log in again after extended inactivity.

A custom RoleMiddleware was implemented to protect routes based
on user roles. This middleware inspects the user's role before
allowing access to specific endpoints. For example, only users
with the admin role can access the administrative dashboard,
while only authenticated clients can create price alerts. This
approach ensures that users cannot accidentally or intentionally
access features intended for other roles.

The User model includes helper methods that make role checking
simple throughout the application. Methods like is admin,
is client, is merchant, and is supplier return true or false
values indicating whether the current user belongs to that
role. These methods are used in controllers, middleware, and
views to conditionally display content and functionality.

Testing validated that the authentication system worked correctly
in all scenarios. New users could register and verify their
phone numbers successfully. Returning users could log in with
valid credentials and receive valid tokens. Users attempting
to access endpoints outside their role were properly rejected
with appropriate error messages. The logout process correctly
invalidated tokens, preventing further use after disconnection.

\begin{figure}[H]
\begin{center}
\fbox{Add screenshot: Login page / Registration form / API response}
\end{center}
\caption{Authentication Interface}
\label{fig:auth_interface}
\end{figure}

% ============================================================
\section*{Sprint 2: Product Catalog and Search}
\addcontentsline{toc}{section}{Sprint 2: Product Catalog and Search}

\subsection*{Sprint Overview}

Sprint 2 built the core functionality that makes PrixTunisix
useful to its visitors: the product catalog and search system.
This sprint focused on creating a comprehensive database of
products that users could browse, filter, and search to find
the items they wanted at the best prices. The catalog serves
as the foundation for all other client-facing features,
including price comparisons, favorites, and alerts.

The product data model needed to support several key requirements.
Products needed to belong to categories and brands for organized
browsing. Each product could have multiple offers from different
merchants at different prices. The system needed to calculate the
lowest price across all available offers for each product.
Price history needed to be tracked over time to show trends.

Categories presented a particular challenge because they needed
to support hierarchical organization. A product category like
Electronics could contain subcategories like Computers, Phones,
and Accessories. Those subcategories could contain further
subcategories. This recursive structure required a self-referential
database design using a parent identifier foreign key.

The search functionality needed to support multiple approaches.
Users should be able to search by typing product names or
descriptions. They should be able to filter by category to see
only products in a specific area. They should be able to filter
by brand to see only products from manufacturers they prefer.
Results should be sortable by price or by newest first.

Several specialized controllers were created to handle different
aspects of the catalog. The ProductController manages product
CRUD operations and validation. The CategoryController handles
the hierarchical category structure. The BrandController manages
manufacturer information. The SearchController processes search
queries and applies filters. Additional controllers handle
boutique pages for specific merchants, marque pages for specific
brands, and an artificial intelligence chatbot for user assistance.

The most important feature implemented was the lowest price
calculation. For every product displayed, the system shows the
cheapest available offer from all merchants. This calculation
is performed by querying the database for the minimum price among
all offers where the offer is marked as available. This simple method
proved essential for the user experience, as users come to the
site specifically to find the best deals.

Price history tracking was added to show how prices have changed
over time. The system stores every observed price for every
offer in an append-only history table. Users can view the last
30 days of price changes as an interactive chart, helping them
decide whether to buy now or wait for a better price.

\begin{table}[H]
\begin{center}
\begin{tabular}{|m{3cm}|m{6cm}|}
\hline
\textbf{Model} & \textbf{Relationships} \\ \hline
Product & belongs to Category, Brand; has many Offers, PriceAlerts, Favorites \\ \hline
Category & belongs to parent Category; has many children Categories, Products \\ \hline
Brand & has many Products \\ \hline
Offer & belongs to Product, Merchant, MerchantWebsite; has many PriceHistory \\ \hline
PriceHistory & belongs to Offer (immutable price records) \\ \hline
\end{tabular}
\caption{Catalog Data Models}
\label{tab:catalog_models}
\end{center}
\end{table}

The search implementation processes queries in several stages.
First, the system determines whether the user is searching by
keyword, filtering by category, or filtering by brand. Then it
constructs the appropriate database query with the relevant
filters applied. For keyword searches, it matches against product
names, descriptions, and reference numbers. For each product
found, the system calculates the lowest available price by
querying all associated offers. Results are returned sorted by
the specified criteria, whether price ascending, price descending,
or newest first.

Pagination ensures that large result sets are delivered in
manageable chunks. Rather than returning thousands of products
at once, the API returns a limited number per page along with
metadata indicating how many total results exist and how many
pages there are. This approach keeps response times fast while
allowing users to browse through all available products.

Testing confirmed that the search system performed well under
load. Keyword searches returned relevant results quickly.
Filters correctly narrowed results to the specified categories
and brands. The lowest price display was accurate across all
products tested. Category navigation worked smoothly with
deep hierarchies. The system handled large product catalogs
without noticeable performance degradation.

\begin{figure}[H]
\begin{center}
\fbox{Add screenshot: Product listing page / Search results / Product details}
\end{center}
\caption{Product Catalog Interface}
\label{fig:catalog_interface}
\end{figure}

% ============================================================
\section*{Sprint 3: Data Scraping System}
\addcontentsline{toc}{section}{Sprint 3: Data Scraping System}

\subsection*{Sprint Overview}

Sprint 3 addressed the challenge of getting pricing data into
the platform automatically. Without an automated scraping system,
the platform would require manual data entry for every merchant
and every product, which would be impossible to maintain at scale.
This sprint built the infrastructure that continuously collects
pricing information from partner merchant websites.

The scraping system works by configuring scripts that define
what websites to scrape, how often to run, and what data to
extract. When a script executes, it fetches the target web pages,
parses the HTML to extract product information and prices, and
saves the results to the database. This happens on a scheduled
basis so that prices stay current without manual intervention.

A critical challenge was matching scraped data to products already
in the catalog. When scraping a merchant website, the system gets
raw product information that may not exactly match how products
have been entered in our system. Product names might be slightly
different, references might use different formats, or products
might be missing entirely. The matching system addresses this by
comparing scraped data against existing products to find the
best matches.

The matching algorithm uses several factors to determine
similarity. Text similarity compares product names to find
likely matches. Reference matching looks for identical product
references across different sources. The system calculates a
confidence score from 0 to 100 percent indicating how sure it
is about each match. Products with high confidence scores are
automatically approved, while lower scores are flagged for human
review.

Suppliers can access the scraping system through an API that
allows external partners to submit their own data. Each supplier
receives an API key for authentication and can configure which
websites their data comes from. This opens the possibility of
working with additional data sources without requiring internal
scraping scripts for every merchant.

\begin{table}[H]
\begin{center}
\begin{tabular}{|m{3.5cm}|m{5.5cm}|}
\hline
\textbf{Model} & \textbf{Description} \\ \hline
MerchantWebsite & Partner website configuration (base URL, logo, active status) \\ \hline
ScrapingScript & Scraping job configuration (target URL, frequency, schedule) \\ \hline
ScrapingLog & Execution history (start time, records collected, errors) \\ \hline
ProductMatch & Match between scraped offer and catalog product with confidence score \\ \hline
Supplier & Extended user with API key for external data submission \\ \hline
\end{tabular}
\caption{Scraping Data Models}
\label{tab:scraping_models}
\end{center}
\end{table}

The scraping execution follows a well-defined process. First,
the system checks that the script is active and not currently
running. It creates a log entry to track the execution. Then it
fetches each target URL and parses the HTML response to extract
product data. The extracted data is compared against existing
products in the catalog using the matching algorithm. New offers
are created for matched products, and low-confidence matches
are flagged for review. Finally, the log is updated with statistics
about how many records were collected and whether any errors
occurred.

The logging system provides visibility into what the scrapers
are doing. Each execution creates a log entry showing when it
started, when it ended, how many products were collected, and
how many errors occurred. If errors happened, the log includes
details about what went wrong. This information is essential
for identifying problems with data sources and ensuring the
scraping pipeline is working correctly.

The frequency setting controls how often each script runs
automatically. Some merchant websites update their prices
frequently and need multiple daily runs, while others change
slowly and might only need weekly checks. The schedule is
configurable for each script, allowing the system to balance
data freshness with server load and respect for partner websites.

By the end of this sprint, the scraping system was fully
functional. Scripts could be created, configured, and executed
either on a schedule or manually. The matching system worked
well enough to handle most products automatically, with human
review catching the difficult cases. The API allowed external
suppliers to connect and submit data. The data pipeline from
merchant websites through to the user-facing catalog was
complete and operational.

Testing verified that scraping scripts ran successfully and
collected data from target websites. The parsing logic handled
different page structures correctly. New offers were created
in the database with proper associations to products and merchants.
The matching system correctly identified obvious matches while
flagging uncertain ones for review. The API responded correctly
to authenticated requests from external suppliers.

\begin{figure}[H]
\begin{center}
\fbox{Add screenshot: Scraping configuration panel / Execution logs}
\end{center}
caption{Scraping System Interface}
\label{fig:scraping_interface}
\end{figure}

% ============================================================
\section*{Conclusion}
\addcontentsline{toc}{section}{Conclusion}
% ============================================================

The first phase of implementation completed three critical
sprints that established the foundation for the entire platform.
Sprint 1 created a secure authentication system capable of
managing five distinct user types with role-based access control.
Sprint 2 built the product catalog and search functionality
that allows visitors to find and compare prices across merchants.
Sprint 3 developed the automated scraping pipeline that keeps
pricing data current without manual data entry.

Together, these three sprints produced a working platform that
could demonstrate core value to users. Visitors could browse
products and see price comparisons. Authentication was in place
to support personalized features in future sprints. The data
collection infrastructure was established to continuously grow
the product database.

The technical decisions made during this phase positioned the
project well for continued development. The role-based permission
system would expand as more features were added. The search
infrastructure would grow more sophisticated as requirements
evolved. The scraping system would add more merchants and data
sources. But the core architecture was solid and ready for the
challenges of subsequent sprints.

With these foundations in place, the project could proceed to
Phase 2 implementation, which would focus on client-facing
features like favorites, alerts, and wishlists, along with
merchant and administrative functionality.